\element{cours}{
    \begin{question}{q01}
        Rappeler ce que signifie qu'une classe est dérivée d'une autre.\\
        \evaluationProfLignes{2}{2}
    \end{question}
}

\element{cours}{
    \begin{question}{q02}
        Représenter graphiquement la structure interne d'un objet de commande.\\
        \evaluationProfLignes{2}{5}
    \end{question}
}

\element{cours}{
    \begin{question}{q03}
        Quand utiliser le pointeur \lstinline|SUPER| ?\\	
        \evaluationProfLignes{2}{2}    
    \end{question}
}

\element{cours}{
    \begin{question}{q03}
        Rappeler le type de POU utilisé pour implémenter une classe dans l'univers CodeSys. Justifier ce choix.\\
        \evaluationProfLignes{2}{2}    
    \end{question}
}

\element{cours}{
    \begin{question}{q04}
        Expliquer la différence entre les accesseurs \lstinline|get| et \lstinline|set|.\\
        \evaluationProfLignes{2}{2}
    \end{question}
}

\element{cours}{
    \begin{question}{q05}
        Déclarer un tableau de 6 booléens. 
        \evaluationProfLignes{2}{2}
    \end{question}
}

\element{cours}{
    \begin{question}{q06}
        Rappeler la définition d'un système temps réel.
        \evaluationProfLignes{2}{2}
    \end{question}
}

\element{cours}{
    \begin{question}{q07}
        Déclarer une structure \lstinline|TVehicule| comportant deux membres : 
        \begin{itemize}
            \item un entier \lstinline|m_nbrDeRoue|, 
            \item un réel \lstinline|m_vitesseMax|.
        \end{itemize}
        \evaluationProfLignes{2}{2}
    \end{question}
}

